% GENERATED LOCAL-BUILD FLAT PART — DO NOT MERGE

% ---- BEGIN TIR/monograph/v12/appendices/appA_formal_proofs_long_algebra.tex ----
% !TEX root = ../../tir_monograph_v12.tex
\chapter{Formal Proofs, Algebraic Firewalls and Legacy Repairs}
\label{app:v12_formal_proofs}

\paragraph{Purpose and provenance.}
This appendix consolidates long algebra that supports Chapters~4--10 and the
diagnostic no-go surfaces used later in the monograph.  Historical appendix
material is preserved through explicit source pointers rather than copied
verbatim.  GREMLIN/HOUND is used here as a bounded adversarial audit layer:
candidate inconsistencies are promoted into this appendix only after an exact
algebraic or deterministic check.

The principal legacy sources are
\path{TIR/monograph/appendices/appB_collatz_tables.tex},
\path{TIR/monograph/appendices/appC_su3_primer.tex},
\path{TIR/monograph/appendices/appD_poincare_disk.tex},
\path{TIR/monograph/appendices/appE_tetrahedron_algebra.tex},
\path{TIR/monograph/appendices/appJ_collatz_quarter_power_scaling.tex},
\path{TIR/monograph/appendices/appK_sector_holonomy_release.tex}, and
\path{TIR/monograph/appendices/appL_up_sector_common_baseline_no_go.tex}.

\section{Finite Collatz statements and the structural labels}

For positive integers define
\[
C(n)=
\begin{cases}
3n+1,&n\ \hbox{odd},\\
n/2,&n\ \hbox{even}.
\end{cases}
\]
The seed required by the canonical v12 label is finite and directly checkable:
\[
3\to10\to5\to16\to8\to4\to2\to1.
\]
There are seven steps before the first occurrence of \(1\), giving the typed
structural label
\[
\boxed{L_3=7.}
\]
The twin-prime source gives
\[
\boxed{L_4=5-3=2,\qquad L_5=5.}
\]
These finite statements are the complete arithmetic input required by the
v12 label layer.  Their publication status is inherited from
Chapter~\ref{ch:v12_discrete_labels}.

A second finite trajectory, historically used to motivate the quark-prime
candidate set, is
\[
7\to22\to11\to34\to17\to52\to26\to13\to40\to20\to10\to5
\to16\to8\to4\to2\to1.
\]
The primes encountered are
\[
7,11,17,13,5,3,
\]
whose sorted set is
\[
\mathcal P_Q=\{3,5,7,11,13,17\}.
\]
This is retained as a finite set-generation observation.  The assignment to
ordered quark slots is owned by the explicit typed map of
Chapter~\ref{ch:v12_discrete_labels}; the exhaustive \(6!=720\) permutation
audit leaves two arithmetic survivors before that ordering rule and one after
it.

\section{Corrected \(SU(3)\) representation algebra}

The three-flavour carrier is \(V_F\cong\mathbb C^3\), with mixing algebra
\(\mathfrak{su}(3)\).  Its real dimension is
\[
\dim_{\mathbb R}\mathfrak{su}(3)=3^2-1=8.
\]
For an \(SU(3)\) irreducible representation with Dynkin labels \((p,q)\),
\[
\boxed{
\dim(p,q)=\frac{(p+1)(q+1)(p+q+2)}{2}.
}
\]
Consequently,
\[
\dim(1,0)=3,\qquad
\dim(0,1)=3,\qquad
\dim(1,1)=8,\qquad
\dim(3,0)=10,\qquad
\dim(2,2)=27.
\]
The mesonic flavour nonet belongs to the reducible decomposition
\[
\boxed{
3\otimes\bar3=8\oplus1,
}
\]
rather than to the \((2,2)\) irrep.  The historical primer row that paired
\((2,2)\) with ``meson nonet'' is therefore retained as a repaired
representation-label mismatch.
\vTwelveStatus{D}{RETROSPECTIVE}{PASS}

For the canonical \(\kappa\) normalization, only the exact carrier facts
\[
N_F=3,\qquad \dim\mathfrak{su}(3)=8,\qquad N_{\rm mix}=3\times8=24
\]
are required.  The full normalization derivation remains owned by
Chapter~\ref{ch:v12_kappa_normalization} and
\path{TIR/foundations/TIR_KAPPA_FLAVOUR_MIXING_NORMALIZATION_V0_1.md}.

\section{Poincar\'e disk normalization repair}

For the unit disk \(\mathbb D=\{z\in\mathbb C:|z|<1\}\) with metric
\[
ds^2=\frac{4\,|dz|^2}{(1-|z|^2)^2},
\]
the curvature is \(K=-1\).  Define the disk cross-ratio modulus
\[
\chi(z_1,z_2)=
\left|
\frac{z_1-z_2}{1-\bar z_1z_2}
\right|.
\]
The geodesic distance compatible with the displayed metric is
\[
\boxed{
d_{\mathbb D}(z_1,z_2)
=2\,\operatorname{artanh}\chi(z_1,z_2).
}
\]
For \(z_1=0\) this reduces to
\[
d_{\mathbb D}(0,r)=\int_0^r\frac{2\,du}{1-u^2}
=2\,\operatorname{artanh}r,
\]
which fixes the normalization directly.  The historical appendix printed the
same expression without the factor \(2\); v12 records that line as a metric
normalization repair.
\vTwelveStatus{D}{RETROSPECTIVE}{PASS}

For a geodesic triangle with angles \(\alpha,\beta,\gamma\),
\[
A_{\mathbb D}=\pi-(\alpha+\beta+\gamma).
\]
A one-angle identification therefore supplies only partial cell data.  v12
assigns any proposed information-area equality to an explicit bridge theorem
with the complete triangle, curvature convention and area normalization stated
before comparison.

\section{Regular tetrahedral Gram algebra}

Let \(\mathbf n_a\in\mathbb R^3\), \(a=1,\ldots,4\), be unit vectors satisfying
\[
\mathbf n_a\cdot\mathbf n_b=-\frac13\qquad(a\ne b).
\]
Their Gram matrix is
\[
G_{ab}=
\begin{cases}
1,&a=b,\\
-1/3,&a\ne b.
\end{cases}
\]
Equivalently,
\[
G=\frac43I_4-\frac13J_4,
\]
where \(J_4\) is the all-ones matrix.  Since \(J_4\) has eigenvalues \(4,0,0,0\),
\[
\operatorname{spec}(G)=\left\{0,\frac43,\frac43,\frac43\right\},
\qquad
\operatorname{rank}G=3.
\]
The null eigenvector is \((1,1,1,1)^T\), hence
\[
\sum_{a=1}^4\mathbf n_a=0.
\]
Every squared edge length is
\[
|\mathbf n_a-\mathbf n_b|^2
=2-2\mathbf n_a\cdot\mathbf n_b
=\frac83,
\]
so the four points form a regular tetrahedron in the three-real-dimensional
carrier.  Permuting the four vertices preserves the Gram matrix, giving
\[
\operatorname{Aut}(\Delta^3)\cong S_4,
\qquad |S_4|=24.
\]
This is the finite-symmetry crosscheck used after the independent
flavour-mixing count \(3\times8=24\).

\section{Legacy operator-tetrahedron firewall}

The historical operator-labelled tetrahedron introduced six sector-specific
edge numbers and then wrote
\[
\cos\theta_{ij}
=
\frac{e_{ik}^2+e_{jk}^2-e_{ij}^2}{2e_{ik}e_{jk}}.
\]
That expression is the planar law of cosines for the angle of a triangle built
from three edge lengths.  A tetrahedral dihedral angle instead requires the
normals of two faces, or an equivalent Cayley--Menger/Gram construction using
the full tetrahedral metric data.  v12 therefore assigns the historical
``dihedral'' line to the diagnostic ledger rather than to the regular
tetrahedral closure theorem.
\vTwelveStatus{D}{RETROSPECTIVE}{QUARANTINED}

The same historical appendix states a volume proportionality to \(\kappa^3\)
without an independently derived Gram determinant and scale map.  v12 assigns
volume/PMNS coupling to the sector-model provenance of Chapters~12--13, while
the geometric tetrahedron theorem is owned by Chapter~6.

\section{Accelerated Collatz quarter-power arithmetic}

For odd \(n\), define the accelerated map
\[
T(n)=\frac{3n+1}{2^{a(n)}},
\qquad
a(n)=\nu_2(3n+1).
\]
Under the standard uniform odd-residue calculation,
\[
\Pr(a=k)=2^{-k},\qquad k\ge1,
\]
so
\[
\mathbb E[a]=\sum_{k\ge1}k2^{-k}=2.
\]
The associated asymptotic geometric-mean multiplier is then
\[
\boxed{
\rho_C
=\exp\!\left(\ln3-\mathbb E[a]\ln2\right)
=\frac34.
}
\]
This arithmetic result is distinct from the particle-model bridge that uses
\(3/4\) as an exponent.  The archived one-anchor charged-fermion trace keeps its
retrospective physical-spectrum failure, while the arithmetic multiplier
remains a technical result.

\section{Common-baseline translation no-go}

Let \(r_g\) be logarithmic residuals for a fixed three-channel up-sector trace,
and apply one common additive baseline \(B\):
\[
r'_g=r_g+B.
\]
For every pair \(g,h\),
\[
r'_g-r'_h=r_g-r_h,
\]
so the residual spread is translation invariant.

\paragraph{Proposition.}
There exists a common \(B\) satisfying
\[
|r_g+B|\le\varepsilon
\quad\text{for every }g
\]
exactly when
\[
\max_g r_g-\min_g r_g\le2\varepsilon.
\]
The smallest achievable uniform error is
\[
\boxed{
\varepsilon_{\min}
=\frac{\max r-\min r}{2}.
}
\]

\paragraph{Proof.}
Each channel requires
\[
B\in[-r_g-\varepsilon,-r_g+\varepsilon].
\]
All intervals intersect exactly when the largest lower endpoint is no greater
than the smallest upper endpoint, which is equivalent to the stated spread
condition.  Centering the two extreme residuals around zero attains half the
spread.

For the archived v10.2 residual vector
\[
(r_u,r_c,r_t)
=(3.1857238256,\,0.0818052322,\,0.2293790515),
\]
the spread is
\[
\Delta r=3.1039185934
\]
and therefore
\[
\varepsilon_{\min}=1.5519592967,
\qquad
e^{\varepsilon_{\min}}=4.7207104\ldots .
\]
The common-baseline architecture is thus excluded for the fixed relative trace
at any target envelope narrower than this minimax bound.
\vTwelveStatus{D}{RETROSPECTIVE}{PASS}

\section{Appendix-A audit ownership}

The deterministic appendix integration gate is
\path{TIR/validation/tir_v12_appendix_integration_audit_v0_1.py}.
Its machine-readable receipt is
\path{TIR/validation/TIR_V12_APPENDIX_INTEGRATION_AUDIT_V0_1.json}.
The gate checks the corrected \(SU(3)\) representation identity, the factor-two
Poincar\'e distance normalization, the regular tetrahedral Gram spectrum, the
legacy dihedral quarantine marker and the common-baseline minimax identity.

% ---- END TIR/monograph/v12/appendices/appA_formal_proofs_long_algebra.tex ----
\clearpage

% ---- BEGIN TIR/monograph/v12/appendices/appB_numerical_tables_reproducibility.tex ----
% !TEX root = ../../tir_monograph_v12.tex
\chapter{Numerical Tables, Receipts and Reproducibility}
\label{app:v12_numerics_reproducibility}

\paragraph{Purpose.}
This appendix is the numerical companion to Chapter~\ref{ch:v12_evidence_matrix}.
Sector chapters own formulas; Chapter~19 owns current evidence verdicts; this
appendix owns reproducible arithmetic, validator locations and frozen receipts.
Historical calculation tables remain available under
\path{TIR/monograph/appendices/}, while v12 records their audited numerical
status rather than reproducing target-matched tables as fresh evidence.

\section{Canonical numerical constants}

The v12 constant surface used by the migrated particle-sector calculations is
\[
L_3=7,\qquad L_4=2,\qquad L_5=5,
\]
and
\[
\boxed{
\kappa=\frac{\ln2}{24\pi}
=0.009193\ldots .
}
\]
The exact ownership chain is
\[
\mathbb C^3
\to SU(3)_F
\to \dim\mathfrak{su}(3)=8
\to N_{\rm mix}=24
\to\Phi_{\rm mix}=24\pi
\to\kappa.
\]
The deterministic source is
\path{TIR/validation/tir_kappa_flavour_mixing_normalization_v0_1.py}.

\section{Active validator index}

The current publication build executes the following v12-specific numerical and
source validators in addition to the foundational geometry suite:
\begin{description}
\item[Discrete labels:]
\path{TIR/validation/tir_v12_discrete_labels_audit_v0_1.py}.
\item[Coefficient forcing:]
\path{TIR/validation/tir_coefficient_role_orientation_forcing_v0_1.py}.
\item[Coefficient selector composition no-go:]
\path{TIR/validation/tir_coefficient_selector_composition_nogo_v0_3.py}.
\item[Neutrino action consistency:]
\path{TIR/validation/tir_v12_neutrino_action_consistency_v0_1.py}.
\item[Neutrino source repair:]
\path{TIR/validation/tir_neutrino_absolute_action_source_repair_v0_1.py}.
\item[CKM/PMNS internal consistency:]
\path{TIR/validation/tir_v12_flavour_mixing_internal_consistency_v0_1.py}.
\item[CKM Jarlskog invariant closure:]
\path{TIR/validation/tir_ckm_jarlskog_invariant_closure_v0_1.py}.
\item[CKM 600-cell incidence:]
\path{TIR/validation/tir_ckm_600cell_flag_incidence_candidate_v0_1.py}.
\item[CKM local A2/operator bridge:]
\path{TIR/validation/tir_ckm_600cell_a2_selector_bridge_candidate_v0_2.py}.
\item[CKM projective even-sector selector:]
\path{TIR/validation/tir_ckm_projective_even_sector_selector_candidate_v0_3.py}.
\item[CKM clean noncommuting family pair:]
\path{TIR/validation/tir_ckm_600cell_clean_noncommuting_family_pair_candidate_v0_4.py}.
\item[CKM lambda0 nested-fraction provenance:]
\path{TIR/validation/tir_ckm_lambda0_nested_fraction_decomposition_v0_5.py}.
\item[CKM half-kernel kappa selector:]
\path{TIR/validation/tir_ckm_half_kernel_kappa_selector_candidate_v0_6.py}.
\item[CKM product-flow uniqueness audit:]
\path{TIR/validation/tir_ckm_half_kernel_product_flow_uniqueness_audit_v0_7.py}.
\item[CKM angle--sine generator firewall:]
\path{TIR/validation/tir_ckm_su3f_angle_sine_generator_firewall_v0_8.py}.
\item[Golden-mean primitive-orbit asymptotic:]
\path{TIR/validation/tir_golden_mean_prime_orbit_asymptotic_v0_1.py}.
\item[Hadron formula audit:]
\path{TIR/validation/tir_v12_hadron_formula_audit_v0_2.py}.
\item[Hypercharge/anomaly audit:]
\path{TIR/validation/tir_v12_hypercharge_anomaly_audit_v0_1.py}.
\item[Hypercharge relative uniqueness:]
\path{TIR/validation/tir_hypercharge_source_uniqueness_v0_1.py}.
\item[Electroweak/Higgs audit:]
\path{TIR/validation/tir_v12_electroweak_higgs_audit_v0_1.py}.
\item[Strong-CP/neutron-EDM audit:]
\path{TIR/validation/tir_v12_strong_cp_nedm_audit_v0_1.py}.
\item[Cosmology formula audit:]
\path{TIR/validation/tir_v12_cosmology_formula_audit_v0_1.py}.
\item[Evidence matrix consistency:]
\path{TIR/validation/tir_v12_evidence_matrix_consistency_v0_2.py}.
\item[Prospective contract:]
\path{TIR/validation/tir_v12_prospective_contract_audit_v0_1.py}.
\item[Source contract:]
\path{TIR/validation/tir_v12_source_contract_v0_7.py}.
\item[Semantic frontier audit:]
\path{TIR/validation/tir_v12_3_semantic_frontier_audit_v0_1.py}.
\item[Appendix integration:]
\path{TIR/validation/tir_v12_appendix_integration_audit_v0_1.py}.
\end{description}

A technical validator PASS certifies its declared arithmetic/source contract.
Observable-level empirical interpretation remains in
Chapter~\ref{ch:v12_evidence_matrix}.

\section{Charged-lepton arithmetic snapshot}

For the historical action coordinates
\[
\begin{aligned}
S_e&=\frac12-3\kappa+\frac{\kappa}{L_3}-\frac{\kappa^2}{2},\\
R_{e\mu}&=5\kappa+\frac{2\kappa}{L_3}
+\frac{L_3+1}{2}\kappa^2,\\
R_{\mu\tau}&=3\kappa-\frac{\kappa}{L_3}
-\frac{L_3}{2}\kappa^2,
\end{aligned}
\]
v12 reproduces
\[
S_e=0.473691600121\ldots,\quad
S_\mu=0.424761179774\ldots,\quad
S_\tau=0.398790835923\ldots .
\]
Under the archived exponential mass map, the corresponding values are
\[
m_e=0.511641997\ldots\ {\rm MeV},
\quad
m_\mu=104.831758\ldots\ {\rm MeV},
\quad
m_\tau=1767.503884\ldots\ {\rm MeV}.
\]
The frozen 2026 evidence audit assigns precision FAIL to all three rows.  The
numbers remain useful as a reproducible historical trace, not as a replacement
for the current evidence verdict.

\section{Neutrino formula/value reconciliation}

The legacy neutrino source defines
\[
S_{\rm bare}=\frac{1+L_4/L_3}{2}=\frac9{14},
\qquad
A_{\rm face}=\frac12\left(\frac{L_4}{L_3}\right)^2=\frac2{49},
\]
and
\[
\delta S=\kappa A_{\rm face}(1-\kappa).
\]
The printed source line uses
\[
S_1^{\rm literal}=S_{\rm bare}+\kappa\,\delta S,
\]
whereas the published mass triplet is reproduced by
\[
\boxed{
S_1^{\rm rec}=S_{\rm bare}+\delta S.
}
\]
The literal rule yields approximately
\[
(0.00521482,\ 0.0104296,\ 0.0521482)\ {\rm eV},
\]
while the reconciled rule yields
\[
(0.00500999,\ 0.01001999,\ 0.05009994)\ {\rm eV}.
\]
The source-repair theorem classifies the extra factor of $\kappa$ as a transcription/composition defect and restores the single-offset rule as the canonical repository derivation. The derivational repair is therefore closed; only physical absolute-mass validation remains open, as recorded in Chapter~21.

\section{Flavour-mixing snapshot}

The migrated PMNS formulas give
\[
\sin^2\theta_{13}=\frac1{49}=0.02040816\ldots,
\]
\[
\sin^2\theta_{12}
=\frac{L_4}{L_3+L_4}
+\left(\frac{L_4}{L_3}\right)^2
=0.30385488\ldots,
\]
\[
\sin^2\theta_{23}
=\frac12+\frac{L_4}{L_3^2}
=0.54081633\ldots,
\]
and
\[
\delta_{\rm PMNS}
=\pi\left[
1+\frac{L_4}{L_3}
+\left(\frac{L_4}{L_3}\right)^2
\right]
=246.122449^\circ\ldots .
\]
The corresponding frozen verdicts are COMPATIBLE, OPEN, TENSION and OPEN in the
order \(\theta_{12},\theta_{23},\theta_{13},\delta_{\rm PMNS}\).

For CKM,
\[
\lambda
=\frac{L_4}{L_3+L_4}
+\frac{L_4}{L_3}\kappa
=0.2248488365\ldots,
\]
\[
|V_{cb}|=\frac2{49}=0.0408163265\ldots,
\qquad
|V_{ub}|=\frac8{2205}=0.0036281179\ldots,
\]
and
\[
\delta_{\rm CKM}=\arccos\frac25=66.4218215^\circ\ldots .
\]
With
\[
s_{12}=\lambda,\qquad s_{23}=|V_{cb}|,\qquad s_{13}=|V_{ub}|,
\qquad c_{ij}=\sqrt{1-s_{ij}^2},
\]
the exact rephasing-invariant Jarlskog quantity of the current CKM matrix is
\[
\boxed{
J_{\rm CP}^{\rm CKM}
=s_{12}s_{23}s_{13}c_{12}c_{23}c_{13}^2\sin\delta_{\rm CKM}
=2.97106576668\times10^{-5}.
}
\]
The earlier independent structural assignment
\[
J_{\rm old}
=\kappa^2\frac{L_4}{L_5}
\left(1-\frac12\left(\frac{L_4}{L_5}\right)^2\right)
=3.110115459\times10^{-5}
\]
is retained only as historical provenance and is deprecated as an independent
Jarlskog invariant.  Its relative gap from the exact current CKM invariant is
\[
4.47088522\%.
\]
The closure validator also checks all nine nontrivial CKM quartets and the
matrix unitarity/determinant conditions.

\section{Hadron and meson reproduction firewall}

The baryon octet calculation is retained with provenance quarantine because
the historical coefficient construction was refined against target masses.
The decuplet trace remains a compatible anchor-linked candidate.  Their
current result status is defined in Chapter~19 rather than by the old appendix
summary percentages.

The legacy pion and kaon exponentials contain an immediate dimensional/arithmetic
witness.  With the displayed Planck-energy scale \(E_P\),
\[
E_Pe^{-6/\pi}
\]
remains of order \(10^{21}\) MeV rather than \(139.57\) MeV.  Likewise,
\[
E_Pe^{-(\pi^2/6-1)}
\]
remains of Planck order rather than \(493.68\) MeV.  The target-valued numbers
printed beside those expressions therefore carry formula FAIL in v12.

The historical heavy/vector meson tables list measured values as model outputs
without a complete executable derivation chain.  v12 keeps those rows
QUARANTINED as reproduction/provenance surfaces pending the meson absolute-action
baseline theorem.

\section{External-data ownership}

The historical
\path{TIR/monograph/appendices/appH_pdg_tables.tex}
is a dated reference snapshot.  The current particle-sector publication
comparison is owned by
\path{TIR/validation/pdg2026/TIR_PDG2026_VALIDATION_MATRIX_V2.md},
frozen 15 August 2026.  Values with renormalization-scheme or scale dependence
retain that typing in the current matrix.

This separation prevents an old numerical table from silently becoming the
current validation source after the theory layer changes.

\section{Exact-head publication reproducibility}

The authoritative v12 workflow is
\path{.github/workflows/compile-tir-monograph-v12.yml}.
For a pull request it checks out the exact PR-head SHA, verifies that checkout,
runs the deterministic validator suite, compiles the PDF and then applies a
publication preflight.

The hardened preflight requires:
\begin{enumerate}
\item zero unresolved references or citations and zero multiply-defined labels;
\item zero Hyperref PDF-string token warnings;
\item zero overfull boxes;
\item a successful \texttt{qpdf --check};
\item embedded fonts with zero Type-3 fonts;
\item correct title and author metadata;
\item an exact-head artifact carrying the actual PR-head SHA.
\end{enumerate}

The appendix integration receipt is
\path{TIR/validation/TIR_V12_APPENDIX_INTEGRATION_AUDIT_V0_1.json}.
Together with the per-sector receipts, it distinguishes formal arithmetic
reproducibility from the empirical verdict layer.

% ---- END TIR/monograph/v12/appendices/appB_numerical_tables_reproducibility.tex ----
\clearpage

% ---- BEGIN TIR/monograph/v12/appendices/appC_publication_protocol_prospective_gates.tex ----
% !TEX root = ../../tir_monograph_v12.tex
\chapter{Publication Protocol and Prospective Gates}
\label{app:v12_publication_protocol}

\paragraph{Purpose and source ownership.}
This appendix gives the operational detail behind
Chapter~\ref{ch:v12_predictions_falsification}.  It consolidates the frozen
v10.7 separable candidate family and the publication protocol into one
preregistration-style contract.  Historical source lineage is retained in
\path{TIR/monograph/appendices/appM_prospective_observable_identifiability.tex},
\path{TIR/monograph/appendices/appN_separable_candidate_family.tex}, and
\path{TIR/monograph/appendices/appO_publication_protocol.tex}.

\section{Prospective contract object}

A prospective test is represented by
\[
\boxed{
\mathcal P=(F,\mathcal O,\mathcal D,\mathcal R,\mathcal V),
}
\]
where \(F\) is the frozen formula/version identity, \(\mathcal O\) the assigned
observable, \(\mathcal D\) the post-freeze data gate, \(\mathcal R\) the decision
rule and \(\mathcal V\) the version/no-refit policy.

The active v12 family was frozen on 29 July 2026.  Its executable source is
\path{TIR/validation/separable_universal_candidate_family_v10_7.py}.
All three candidates carry
\vTwelveStatus{E}{PROSPECTIVE}{OPEN}
until the assigned qualifying post-freeze likelihood is evaluated.

\section{Shared separable architecture}

Every candidate has the same typed architecture
\[
\boxed{
\ln y_{f,g}=F(S_f)+D(G_g,R_g).
}
\]
The sector coordinates are
\[
S_\ell=4.5,\qquad
S_d=3.0876164691516155,\qquad
S_u=2.920949802484949,
\]
the generation coordinates are
\[
G_1=0.8471633855025916,\quad
G_2=0.6088811663290957,\quad
G_3=0.23926393244694075,
\]
and the Ramanujan release coordinates are
\[
R_1=0.4754443238393299,\quad
R_2=1.0262897039931138,\quad
R_3=2.442442115365772.
\]
The same functions \(F\) and \(D\) act across charged leptons, down-type quarks
and up-type quarks.

\section{Exactly three frozen candidates}

\paragraph{C1: linear separable action.}
\[
F_1(S)=-S,
\qquad
D_1(G,R)=-G+R.
\]

\paragraph{C2: Collatz quarter-power action.}
\[
F_2(S)=-S^{3/4},
\qquad
D_2(G,R)=-G^{3/4}+R.
\]

\paragraph{C3: inverse quarter-power information potential.}
\[
F_3(S)=\frac{S^{-3/4}}{L_3\kappa},
\qquad
D_3(G,R)=\frac{G^{-3/4}}{L_3\kappa}+R.
\]

The family size is itself frozen:
\[
\boxed{N_{\rm candidate}=3.}
\]
A later functional form receives a new version identity and a new prospective
contract.

\section{Orthogonal observable pair}

Charm and the muon share the second-generation coordinate.  Therefore
\[
\begin{aligned}
\ln y_c&=F(S_u)+D(G_2,R_2),\\
\ln y_\mu&=F(S_\ell)+D(G_2,R_2),
\end{aligned}
\]
so
\[
\boxed{
\ln\frac{y_c}{y_\mu}=F(S_u)-F(S_\ell).
}
\]
This observable isolates the sector contribution.

Charm and top share the up-type sector coordinate:
\[
\begin{aligned}
\ln y_c&=F(S_u)+D(G_2,R_2),\\
\ln y_t&=F(S_u)+D(G_3,R_3),
\end{aligned}
\]
hence
\[
\boxed{
\ln\frac{y_c}{y_t}
=D(G_2,R_2)-D(G_3,R_3).
}
\]
This observable isolates the generation/release contribution.

The active class-A mapping is therefore \(y_c/y_\mu\), while \(y_c/y_t\) is the
generation/release test.  The earlier \(y_c/y_\tau\) mapping remains historical
provenance from the pre-orthogonalization stage.

\section{Frozen numerical predictions}

\begin{center}
\begin{tabular}{lrr}
\toprule
Candidate & \(y_c/y_\mu\) & \(y_c/y_t\)\\
\midrule
C1 & 4.850346751338371 & 0.16766796647328305\\
C2 & 2.3521800134268784 & 0.17147213462587316\\
C3 & 6.858021228826228 & \(2.8101955040512466\times10^{-11}\)\\
\bottomrule
\end{tabular}
\end{center}

The deterministic v12 prospective validator reconstructs all six values from
the frozen coordinates and candidate definitions.

\section{Dataset admission gate}

For each observable, the admissible dataset is the first qualifying
post-29-July-2026 joint ATLAS/CMS direct likelihood that reports the relevant
coupling ratio or sufficient likelihood information to construct it on a common
scheme/scale surface.

Dataset admission is decided before reading the numerical location of the
likelihood maximum relative to the three predictions.  If a release lacks the
required joint information, the contract remains OPEN for that release.

\section{Decision rule}

For candidate \(C_i\) and assigned observable \(O\),
\[
\boxed{
\operatorname{PASS}(C_i,O)
\iff
r_i(O)\in{\rm CR}_{95\%}(O),
}
\]
where \(r_i(O)\) is the frozen predicted ratio and
\({\rm CR}_{95\%}(O)\) is the published 95\% confidence region of the admitted
likelihood.

A candidate outside the admitted confidence region receives FAIL for that
frozen test.  The formula identity remains fixed in the corresponding receipt.

\section{Multiplicity-aware interpretation}

There are three simultaneously frozen candidates and two orthogonal observables.
The publication report therefore presents the full \(3\times2\) result surface,
rather than selecting only the most favourable surviving candidate after the
data are visible.

Candidate ranking, if performed, must use the same admitted likelihood surface
for all candidates and must preserve the original three-member family in the
reported ledger.

\section{No-refit and versioning rule}

After freeze:
\begin{enumerate}
\item \(F_i\), \(D_i\), \(S_f\), \(G_g\), \(R_g\), \(L_3\) and \(\kappa\) remain fixed for this contract;
\item the observable identities \(y_c/y_\mu\) and \(y_c/y_t\) remain fixed;
\item an empirical FAIL is recorded without coefficient tuning;
\item a new formula, coordinate or observable receives a new version identity;
\item already inspected retrospective mass tables remain outside the prospective score.
\end{enumerate}

This creates an append-only evidence history: a failed frozen candidate can
coexist with a later new candidate version without rewriting the earlier test.

\section{External-input and scheme discipline}

Yukawa couplings and quark quantities are scheme- and scale-sensitive.
The qualifying experimental surface must therefore provide a common likelihood
or a declared translation to a common scheme before a precision decision is
made.  Ratios constructed from incompatible conventions remain outside the
frozen decision gate.

Upper limits, circular phases and dimensionful anchor quantities retain their
own statistical typing and are not pooled into one global percentage score.

\section{Machine-readable gate}

The prospective contract is checked by
\path{TIR/validation/tir_v12_prospective_contract_audit_v0_1.py}.
The appendix integration audit additionally checks that:
\begin{itemize}
\item exactly three candidate definitions are present;
\item the 29 July 2026 freeze date is present;
\item both orthogonal observables are present;
\item all six frozen ratios are present;
\item the active sector observable is \(y_c/y_\mu\);
\item the no-refit/version rule is explicit.
\end{itemize}

The appendix-level receipt is
\path{TIR/validation/TIR_V12_APPENDIX_INTEGRATION_AUDIT_V0_1.json}.

% ---- END TIR/monograph/v12/appendices/appC_publication_protocol_prospective_gates.tex ----
\clearpage

% ---- BEGIN TIR/monograph/v12/appendices/appD_cross_framework_interfaces.tex ----
% !TEX root = ../../tir_monograph_v12.tex
\chapter{Cross-Framework Interfaces and Dependency Direction}
\label{app:v12_cross_framework}

\paragraph{Purpose.}
This appendix records typed interfaces between TIR and sibling analytic or
dynamical frameworks.  Each interface preserves source direction: exact
mathematical identities remain exact, TIR structural identifications retain
their TIR ownership, and downstream physical bindings keep their own validation
gate.  GREMLIN may propose cross-domain isomorphisms on this surface; theorem or
validator receipts remain the promotion gate.

\section{Half coordinate and binary information}

For binary entropy
\[
H_2(p)=-p\ln p-(1-p)\ln(1-p),
\]
\[
H_2'(p)=\ln\frac{1-p}{p},
\qquad
H_2''(p)=-\frac1p-\frac1{1-p}<0.
\]
Therefore
\[
\boxed{
\operatorname*{arg\,max}_{0<p<1}H_2(p)=\frac12,
\qquad
H_2(1/2)=\ln2.
}
\]
The same midpoint is the unique fixed point of complement \(p\mapsto1-p\).
In the odds coordinate
\[
q=\frac{p}{1-p},
\]
it maps to \(q=1\), the positive fixed point of reciprocal inversion.

TIR uses this exact half/information pair as the numerator side of the canonical
flavour-mixing normalization.

\section{TIR ownership of the \(24\pi\) denominator}

The flavour carrier is
\[
V_F\cong\mathbb C^3,
\qquad
U_F\in SU(3)_F.
\]
Hence
\[
N_F=3,\qquad
\dim_{\mathbb R}\mathfrak{su}(3)_F=8.
\]
The generator--flavour incidence count is
\[
N_{\rm mix}=3\times8=24.
\]
The half coordinate on a full \(2\pi\) angular closure supplies the primitive
half-turn
\[
\Delta\phi_{1/2}=\pi,
\]
so
\[
\Phi_{\rm mix}=24\pi.
\]
Combining the balanced binary information with the mixing-phase measure gives
\[
\boxed{
\kappa
=\frac{H_2(1/2)}{\Phi_{\rm mix}}
=\frac{\ln2}{24\pi}.
}
\]
This is the canonical TIR-internal normalization theorem.  The regular
tetrahedral symmetry \(|S_4|=24\) supplies an independent finite-symmetry
crosscheck of the same integer while Chapter~9 retains sole publication
ownership of the complete derivation.

\section{Secret-of-a-Half negative-inverse interface}

Define
\[
\Omega(s)=\frac{s}{1-s},
\qquad
z_L(s)=1-\frac1s=-\frac1{\Omega(s)}.
\]
Writing \(s=\sigma+it\),
\[
|z_L(s)|^2
=
\frac{(\sigma-1)^2+t^2}{\sigma^2+t^2}.
\]
Therefore
\[
|z_L(s)|=1
\iff
(\sigma-1)^2=\sigma^2
\iff
\boxed{\sigma=\frac12}.
\]
Thus the critical axis is exactly the unit-circle locus of the
negative-inverse coordinate:
\[
\boxed{
\Re s=\frac12
\iff
|z_L(s)|=1.
}
\]

The current cross-framework source is
\path{TIR/interfaces/TIR_SOH_GLOBAL_LI_DOMINATION_BRIDGE_V0_1.md}.
For one reciprocal-conjugate zero quartet with
\[
z_\rho=Re^{i\phi},
\]
its local Li contribution is
\[
L_n(Q_\rho)
=4-2(R^n+R^{-n})\cos(n\phi).
\]
The interface document gives a conditional global-domination theorem candidate:
under the standard symmetric Li representation and classical zero-counting
estimate, an off-axis zero creates an extremal radial shell whose recurrent
phase subsequence dominates the subextremal remainder and yields negative Li
coefficients along a subsequence.

The remaining analytic promotion target in that branch is the independently
derived native positivity implication
\[
\boxed{
\text{native arithmetic closure}
\Longrightarrow
\lambda_n\ge0\quad\forall n.
}
\]
This target remains OPEN in Chapter~21.
\vTwelveStatus{B}{--}{OPEN}

\section{Self-dual fixed point versus extremality}

The half-axis interface also carries a useful diagnostic boundary.  Reciprocal
symmetry by itself fixes a self-dual locus, while dynamical preference requires
an additional variational or positivity condition.

In a finite reciprocal family \(N_n(q)=N_n(1/q)\), the fixed point \(q=1\) may
coexist with off-centre reciprocal extrema.  v12 therefore treats
\[
\boxed{
\text{reciprocal symmetry}
+\text{explicit extremality condition}
\to
\text{extremum claim}
}
\]
as the admissible dependency order.  This prevents symmetry from silently
serving as an unproved stability theorem.

\section{TIR--IDT phase-clock scale interface}

The projective quantities
\[
a_{FS},\qquad
\int\mathcal F_B,\qquad
|\langle\psi_a|\psi_b\rangle|^2,\qquad
\kappa
\]
are dimensionless.  The IDT phase-clock interface supplies the length carrier
\[
\boxed{
\ell_\varphi
=\frac{c}{|\omega_t|}
=c\left|\frac{dt}{d\varphi}\right|
=\frac{\hbar c}{E}.
}
\]
The exact dimensional typing therefore permits the physicalization candidate
\[
\boxed{
ds_{\rm rel}^2
=\ell_\varphi^2\,ds_{FS}^2.
}
\]
For
\[
ds_{FS}^2
=\frac14(d\theta^2+\sin^2\theta\,d\varphi^2),
\]
the associated area candidate is
\[
d\mathcal A_{\rm rel}
=\ell_\varphi^2\,da_{FS}.
\]
With spin-\(1/2\) Berry curvature
\[
\mathcal F_B=\pm2\,da_{FS},
\]
one obtains
\[
d\mathcal A_{\rm rel}
=\pm\frac{\ell_\varphi^2}{2}\mathcal F_B.
\]

The source contract is
\path{TIR/integration/TIR_IDT_INFORMATION_AREA_CURVATURE_V0_2.md},
whose status separates exact dimensional closure from the physical
phase-clock scale-binding candidate.

\section{Polyhedral refinement export}

For a CP1 polyhedral cell complex \(P\) with facewise phase-clock scale,
\[
\mathcal A_{\rm rel}^{(P)}
=
\sum_{f\in P}\ell_{\varphi,f}^2\,a_{FS}(f).
\]
For a uniform cell rate,
\[
\mathcal A_{\rm rel}^{(P)}
=
\frac{c^2}{\omega_P^2}a_{FS}^{(P)}.
\]
The refinement target is
\[
\boxed{
\mathcal A_{\rm rel}^{(P_n)}
\longrightarrow
\mathcal A_{\rm rel}
}
\]
under a controlled refinement sequence.  This target is the dimensional bridge
feeding the TIR--IDT spatial/temporal join.

\section{GR/RFC dependency export}

The v12 completion DAG keeps the relativistic line ordered as
\[
\boxed{
\begin{aligned}
\text{discrete solder/torsion}
&\to\text{Cartan refinement}\\
&\to T^a=0\text{ spatial sector}\\
&\to\text{TIR--IDT ADM join}\\
&\to\text{Einstein constraint/evolution closure}.
\end{aligned}
}
\]
The cross-framework appendix therefore exports typed parent surfaces rather
than importing downstream gravitational equations as foundational premises.

The information-area curvature candidate
\[
\Xi_I^{(P)}
=
\frac{\mathcal J_\pi}{a_{FS}^{(P)}}
\left(\frac{\omega_P}{c}\right)^2
\]
is available as a dimensionally closed interface.  RFC owns the downstream
binding of this quantity to its dynamical curvature/action surface.

\section{Software and GREMLIN boundary}

The runtime/solver layer may consume the normalized cycle, phase and relation
surfaces exported by TIR.  Runtime behavior is not a proof parent for the
mathematical theorem graph; instrument or runtime tests enter through explicit
evidence receipts.

GREMLIN is used as a bounded candidate-generation and adversarial-audit layer:
SPIDER may propose dependency isomorphisms, HOUND may search for contradictions,
MOLE may propose derivation routes and MANTIS may detect duplicate ownership.
A GREMLIN candidate enters the v12 theorem graph only after an independent
deterministic proof/validator gate.

\section{Interface status table}

\begin{center}
\begin{tabular}{p{0.32\textwidth}p{0.23\textwidth}p{0.32\textwidth}}
\toprule
Interface & Status & Owner\\
\midrule
\(1/2\leftrightarrow\ln2\) & exact & binary information layer\\
\(3\times8\times\pi\to24\pi\) & TIR structural PASS & Chapter~9\\
\(\Re s=1/2\leftrightarrow|z_L|=1\) & exact & SOH interface\\
off-axis \(\to\) negative Li subsequence & conditional theorem candidate & SOH bridge\\
native closure \(\to\lambda_n\ge0\) & OPEN & SOH completion frontier\\
\(\ell_\varphi=c/|\omega_t|\) typing & exact dimensional closure & IDT interface\\
\(ds_{\rm rel}^2=\ell_\varphi^2ds_{FS}^2\) & scale-binding candidate & TIR--IDT interface\\
Cartan/ADM/Einstein closure & OPEN & Chapter~21 / RFC join\\
\bottomrule
\end{tabular}
\end{center}

The deterministic appendix integration validator checks the exact half-axis
identity, canonical \(\kappa\) ownership markers, the IDT length-scale equation,
the Li-positivity frontier marker and the bounded GREMLIN promotion rule.

% ---- END TIR/monograph/v12/appendices/appD_cross_framework_interfaces.tex ----
\clearpage
